% Perron-Frobenius for positive kernels, and the vacuum at every spatial extent.
% Lane: O (operator bridge), modules YangMills/OS/**.
% REGIME: tricotomy; NO scores, NO desk language, NO commit chronology beyond
% the reproducibility freeze.
\documentclass[11pt]{article}
\usepackage[margin=1.1in]{geometry}
\usepackage{amsmath,amssymb,amsthm}
\usepackage{booktabs,longtable,array}
\usepackage[colorlinks=true,linkcolor=blue,urlcolor=blue,citecolor=blue]{hyperref}

\newtheorem{theorem}{Theorem}[section]
\newtheorem{lemma}[theorem]{Lemma}
\newtheorem{proposition}[theorem]{Proposition}
\newtheorem{corollary}[theorem]{Corollary}
\newtheorem{remark}[theorem]{Remark}
\theoremstyle{definition}
\newtheorem{definition}[theorem]{Definition}

\emergencystretch=3em
\tolerance=2000

\newcommand{\lean}[1]{\texttt{#1}}
\newcommand{\anchor}{6f3121d693e4f0f560a2ff01c63f4aacf1a49390}
\newcommand{\osline}[3]{%
  \href{https://github.com/lluiseriksson/THE-ERIKSSON-PROGRAMME/blob/\anchor/YangMills/OS/#1.lean\#L#2}{\nolinkurl{#3}}}
\newcommand{\repofile}[2]{%
  \href{https://github.com/lluiseriksson/THE-ERIKSSON-PROGRAMME/blob/\anchor/#1}{\nolinkurl{#2}}}
\newcommand{\R}{\mathbb{R}}
\newcommand{\C}{\mathbb{C}}
\newcommand{\Z}{\mathbb{Z}}
\newcommand{\SU}{\mathrm{SU}}
\newcommand{\Om}{\Omega}

\title{The Vacuum Was Never Absent:\\
Perron--Frobenius for Strictly Positive Kernels in Lean~4,\\
and the Coupled-Slice Vacuum at Every Spatial Extent\\
\large existence by compactness, in a library with no fixed-point theorem}
\author{Lluis Eriksson}
\date{}

\begin{document}
\maketitle

\begin{abstract}
Two companion papers established, for a $\Z_2$ lattice gauge slice with a
spatial coupling, that the elementary route to the vacuum stops
\cite{ErikssonSpatial} and that the natural replacement --- the Hilbert
projective metric --- is blind to the coupling and degenerates in the volume
\cite{ErikssonBirkhoff}.  Both had to work around the same absence: the pinned
\texttt{mathlib} carries no Perron--Frobenius theorem.  The first paper could
therefore only say the vacuum had become unavailable; the second had to build
its domination bound from scratch, and could exhibit the vacuum in closed form
only at two sites.

This paper discharges that dependency.  For a strictly positive kernel on a
finite nonempty type we prove, in Lean~4 with \texttt{mathlib}: a strictly
positive eigenvector \emph{exists}; its eigenvalue is strictly positive; any two
strictly positive eigenvectors are proportional and share their eigenvalue; and
every real eigenvector for that eigenvalue is a scalar multiple of it.  Together
with the domination theorem of \cite{ErikssonBirkhoff} this gives the
Perron statement the lane needs: the eigenvalue is the spectral radius.

The existence proof does not use a fixed-point theorem, because the library has
none.  It maximises $r$ over the compact set of pairs $(r,x)$ with $x$ in the
simplex and $rx\le Ax$; maximality forces equality, since a strict inequality
anywhere would let one further application of $A$ produce an admissible pair with
a larger $r$.  The bound that keeps the set compact is obtained by summing the
constraint: $r=r\sum x\le\sum(Ax)$.

The application is the point.  At \emph{every} spatial extent, and for
\emph{every} strictly positive weight on the source configuration --- the class
that contains the coupled kernel of \cite{ErikssonSpatial} --- the vacuum exists,
is unique up to scale, and carries the spectral radius.  The obstruction of the
first paper was never an absence; it was an unavailability, and it was an
unavailability of one route rather than of the object.

\textbf{No spectral gap is proved here, uniform in the volume or otherwise, and
none is claimed.}  Nothing in this paper is a claim about $\SU(N)$, the
continuum limit, or the Yang--Mills mass gap.
\end{abstract}

\section{Introduction}

\subsection{A dependency that bit twice}

\cite{ErikssonSpatial} proved that switching on a coupling inside a slice
destroys the constancy of the row sums, so the uniform vector stops being fixed
and the elementary construction of the vacuum halts.  It phrased the conclusion
carefully: the vacuum becomes a Perron vector that row-sum normalisation no
longer supplies.  It did not claim the vacuum was gone --- but it could not
produce it either.

\cite{ErikssonBirkhoff} then showed the natural replacement fails for two
independent reasons, and wrote the vacuum down at $L=2$ in closed form.  To
identify its eigenvalue as the spectral radius it had to prove a domination
bound from first principles, for the same reason: no Perron--Frobenius in the
library.

A dependency that forces a detour twice is not bad luck.  This paper removes it.

\subsection{What is proved}

Throughout, $\iota$ is a finite nonempty type and $A:\iota\times\iota\to\R$ is
\emph{strictly positive}, $A_{ij}>0$.  Vectors are functions $\iota\to\R$ and
$(Ax)_i=\sum_j A_{ij}x_j$.

\begin{itemize}
\item \textbf{Existence.}  There are $v$ and $\lambda$ with $v_i>0$ for all $i$,
      $\lambda>0$, and $Av=\lambda v$.
\item \textbf{Uniqueness up to scale.}  If $v,w$ are strictly positive
      eigenvectors with eigenvalues $\lambda,\mu$, then $\lambda=\mu$ and
      $w=cv$ for some $c>0$.
\item \textbf{Geometric simplicity.}  Every real eigenvector for $\lambda$ is a
      scalar multiple of $v$.
\item \textbf{Spectral radius.}  Every eigenvalue, real or complex, is dominated
      in modulus by $\lambda$ \cite{ErikssonBirkhoff}.
\end{itemize}

\medskip
\noindent\textbf{The application.}  For every spatial extent $L$ and every
strictly positive $w$ on configurations, the kernel
$K_w(\sigma,\tau)=w(\sigma)K(\sigma,\tau)$ --- with $K$ the decoupled kernel of
\cite{ErikssonSpatial} --- has all four properties.  The coupled kernel of that
paper is the case of a single spatial bond, and a bond around the slice is
another.

\subsection{What is deliberately not claimed}

\begin{enumerate}
\item \textbf{No spectral gap.}  Nothing here bounds the subdominant eigenvalue,
at any $L$, in any region of parameters.  That is a different problem and this
paper does not touch it.  In particular nothing here contradicts, or improves
on, the obstructions of \cite{ErikssonBirkhoff}.
\item \textbf{Algebraic simplicity is not proved.}  The eigenspace is shown to
be one-dimensional; that the eigenvalue is a simple root of the characteristic
polynomial is not formalised.
\item \textbf{Strict positivity is assumed, not irreducibility.}  The
Perron--Frobenius theory of merely irreducible nonnegative kernels, and the
periodicity theory, are outside this development.
\item \textbf{Still $\Z_2$}, one variable per site, in the application.
\item \textbf{Nothing about $\SU(N)$, the continuum limit, or the Clay problem}
\cite{Clay}.  No reduction is claimed and none is known to the author.
\end{enumerate}

\section{One comparison, used three times}

\begin{lemma}[Positivity of the image]\label{lem:img}
If $z\ge0$ and $z_{j_0}>0$ for some $j_0$, then $(Az)_i>0$ for every $i$.
\end{lemma}

\osline{PerronKernel}{78}{apply\_pos\_of\_nonneg}.

\begin{lemma}[The vanishing comparison]\label{lem:van}
If $z\ge0$, $Az=\lambda z$ and $z_{i_0}=0$ for some $i_0$, then $z=0$.
\end{lemma}

\begin{proof}
Otherwise $z$ is positive somewhere, so $(Az)_{i_0}>0$ by
Lemma~\ref{lem:img}; but $(Az)_{i_0}=\lambda z_{i_0}=0$.
\end{proof}

\osline{PerronKernel}{100}{eq\_zero\_of\_nonneg\_of\_eigen\_of\_zero}.  This
single lemma carries both uniqueness statements: in each case one subtracts the
largest multiple of one eigenvector that still fits under the other, which
produces a nonnegative eigenvector with a zero entry.

\begin{lemma}[Positivity of the eigenvalue]\label{lem:eigpos}
If $v>0$ and $Av=\lambda v$ then $\lambda>0$.
\end{lemma}

\osline{PerronKernel}{87}{eigenvalue\_pos}.

\section{Uniqueness and geometric simplicity}

\begin{theorem}[Geometric simplicity]\label{thm:simple}
If $v>0$ with $Av=\lambda v$, then every $w$ with $Aw=\lambda w$ satisfies
$w=cv$ for some $c\in\R$.
\end{theorem}

\begin{proof}
Let $c=\min_i w_i/v_i$, attained at $i_0$.  Then $z=w-cv\ge0$ by the choice of
$c$, $z_{i_0}=0$, and $Az=\lambda z$ by linearity.  Lemma~\ref{lem:van} gives
$z=0$.
\end{proof}

\osline{PerronKernel}{113}{eigenvector\_eq\_smul\_of\_pos}.

\begin{theorem}[Uniqueness up to scale]\label{thm:unique}
Two strictly positive eigenvectors have the same eigenvalue, and are positive
multiples of one another.
\end{theorem}

\begin{proof}
Both eigenvalues are positive by Lemma~\ref{lem:eigpos}, and each dominates the
other in modulus by the domination theorem of \cite{ErikssonBirkhoff}, so they
are equal; then Theorem~\ref{thm:simple} applies, and the constant is positive
because both vectors are.
\end{proof}

\osline{PerronKernel}{148}{pos\_eigenvector\_unique}.

\section{Existence, by compactness}

The library has no Brouwer fixed-point theorem, so existence is obtained without
one.

\begin{definition}
$\Delta=\{x:\iota\to\R\mid x\ge0,\ \sum_i x_i=1\}$.
\end{definition}

\osline{PerronKernel}{176}{simplexSet}, compact by
\osline{PerronKernel}{184}{isCompact\_simplexSet}: it is a closed subset
of $[0,1]^{\iota}$, and $x_i\le\sum_j x_j=1$ puts it there.

\begin{lemma}[Summing the constraint bounds the multiplier]\label{lem:bound}
If $x\in\Delta$ and $rx\le Ax$ then $r\le R:=\sum_{i}\beta$, where $\beta$ is the
largest entry of $A$.
\end{lemma}

\begin{proof}
Sum the constraint over $i$: the left side is $r\sum_i x_i=r$, and the right side
is $\sum_i(Ax)_i\le\sum_i\beta\sum_j x_j=R$.
\end{proof}

\osline{PerronKernel}{201}{exists\_pos\_eigenvector} contains this step.  It
is what makes the constraint set compact without any estimate on how small the
entries of $x$ can be --- the usual device of restricting to a thickened simplex
is not needed.

\begin{theorem}[Existence]\label{thm:exist}
A strictly positive kernel has a strictly positive eigenvector, with strictly
positive eigenvalue.
\end{theorem}

\begin{proof}
Let $K=\{(r,x)\in[0,R]\times\Delta\mid rx\le Ax\}$.  It is nonempty (take $r=0$
and $x$ a vertex) and compact (a closed constraint on a product of compacta), so
$r$ attains a maximum on it, at $(\lambda,v)$ say.

Since $v\in\Delta$ it is nonnegative and not zero, so $Av>0$ entrywise by
Lemma~\ref{lem:img}.  Suppose $Av\neq\lambda v$.  Then $z=Av-\lambda v\ge0$ is
nonzero, so $Az>0$ entrywise, i.e.\ $A(Av)>\lambda\,Av$ entrywise.  Put
$y=Av/\sum_i(Av)_i\in\Delta$; by homogeneity $Ay>\lambda y$ entrywise, and $y>0$,
so
$\varepsilon=\min_i\bigl((Ay)_i-\lambda y_i\bigr)/y_i>0$
satisfies $(\lambda+\varepsilon)y\le Ay$.  By Lemma~\ref{lem:bound},
$\lambda+\varepsilon\le R$, so $(\lambda+\varepsilon,y)\in K$ --- contradicting
maximality.  Hence $Av=\lambda v$, and then $\lambda v_i=(Av)_i>0$ forces
$\lambda>0$ and $v>0$.
\end{proof}

\osline{PerronKernel}{201}{exists\_pos\_eigenvector}.

\begin{remark}[Why the maximiser is automatically positive]
The simplex admits points with zero entries, and nothing in the definition of
$K$ excludes them.  They are excluded \emph{a posteriori}: once $Av=\lambda v$
is known, $v_i=(Av)_i/\lambda>0$.  Restricting the search to strictly positive
vectors in advance would have cost a compactness argument that this order of
steps avoids.
\end{remark}

\section{The vacuum at every spatial extent}

\begin{definition}
For a weight $w$ on configurations, $K_w(\sigma,\tau)=w(\sigma)K(\sigma,\tau)$
with $K$ the decoupled kernel of \cite{ErikssonSpatial}.
\end{definition}

\osline{PerronKernel}{357}{sourceWeightedKernelL}, strictly positive when $w$
is: \osline{PerronKernel}{361}{sourceWeightedKernelL\_pos}.

\begin{theorem}[The vacuum, at every $L$]\label{thm:vac}
For every $L$, every $\beta$, and every strictly positive $w$, the kernel $K_w$
has a strictly positive eigenvector with strictly positive eigenvalue; it is
unique up to positive scale; and its eigenvalue dominates every eigenvalue, real
or complex.
\end{theorem}

\osline{PerronKernel}{371}{vacuum\_exists\_of\_sourceWeight},
\osline{PerronKernel}{378}{vacuum\_unique\_of\_sourceWeight},
\osline{PerronKernel}{389}{vacuum\_spectral\_radius\_of\_sourceWeight}.
For the concrete slice with a bond around the ring,
\osline{PerronKernel}{401}{spatialWeightRing} and
\osline{PerronKernel}{410}{coupled\_ring\_vacuum\_exists}.

\begin{remark}[What this settles about the first paper]
\cite{ErikssonSpatial} proved that the elementary route stops producing the
vacuum, and was careful to claim no more: it did not assert the vacuum was
absent, and it said a positive kernel on a finite set has a Perron vector.  That
sentence was prose there, cited from the classical theory.  Here it is a
theorem, and it applies to exactly the kernels that paper studied, at every $L$.
So the obstruction is now located precisely: not in the object, and not in the
model, but in the route.
\end{remark}

\begin{remark}[And what it does not settle]
Nothing in this paper produces a spectral gap.  The obstructions of
\cite{ErikssonBirkhoff} are untouched: the projective metric is still blind to
the coupling and still degenerates in the volume, and knowing that the vacuum
exists does not repair either failure.  Having the first object of the chain
back is a precondition for the analytic work, not a substitute for it.
\end{remark}

\section{Formalisation notes}

\paragraph{The order of the steps did the work.}  The compactness argument is
run on the plain simplex, which contains points with zero entries, and strict
positivity of the maximiser is recovered afterwards from the eigen-equation.
Attempting the reverse --- searching only among strictly positive vectors ---
requires a thickened simplex and an invariance estimate; the version here needs
neither.

\paragraph{Summing the constraint replaced an estimate.}  The bound keeping the
constraint set compact is $r=r\sum_i x_i\le\sum_i(Ax)_i$, obtained by summing
the very inequality that defines the set.  A first design bounded $r$ through
the smallest entry of a thickened simplex and needed the ratio of the extreme
entries of $A$; summing removed that entire layer, together with two auxiliary
constants.

\paragraph{One lemma, both uniqueness statements.}  Lemma~\ref{lem:van} is used
for geometric simplicity and, through it, for uniqueness up to scale.  In both
cases the input is produced the same way: subtract the largest multiple that
still fits underneath, and read the resulting zero entry.

\paragraph{A dependency discharged rather than routed around.}  The two
companion developments each paid for the absence of Perron--Frobenius in a
different currency --- a weaker statement in one case, a bespoke proof in the
other.  Building the theorem once cost less than the second detour did.

\section{Reproducibility}\label{sec:repro}

All artifacts are at source checkpoint \texttt{\anchor} on branch \texttt{main}
of \url{https://github.com/lluiseriksson/THE-ERIKSSON-PROGRAMME}.  Every link
was resolved against that commit, and checked to point at an actual declaration
line, before compilation.

\begin{center}
\begin{tabular}{ll}
\toprule
Quantity & Value \\
\midrule
Full core build & 8427 jobs, success \\
Oracle commands & 2559 \\
\quad with axiom dependencies & 2536 \\
\quad axiom-free & 23 \\
Distinct axioms across all outputs & \texttt{propext}, \texttt{Quot.sound}, \texttt{Classical.choice} \\
Occurrences of \texttt{sorryAx} & 0 \\
Project axioms & 0 \\
Errors & 0 \\
\bottomrule
\end{tabular}
\end{center}

Verdicts are in
\repofile{docs/VERIFICATION-LEDGER.md}{docs/VERIFICATION-LEDGER.md}.  The oracle
script is \repofile{oracle\_check.lean}{oracle\_check.lean}.

\begin{thebibliography}{9}

\bibitem{ErikssonSpatial}
L.~Eriksson,
\emph{Where the Elementary Reconstruction Stops: Spatial Coupling Breaks the
Uniform Vacuum, Machine-Checked}, viXra, 2026.

\bibitem{ErikssonBirkhoff}
L.~Eriksson,
\emph{Blind to the Coupling: a Second Machine-Checked Obstruction at Spatial
Extent}, viXra, 2026.

\bibitem{PerronFrobenius}
E.~Seneta,
\emph{Non-negative Matrices and Markov Chains},
2nd ed., Springer, 2006.

\bibitem{Lean4}
L.~de~Moura and S.~Ullrich,
\emph{The Lean 4 theorem prover and programming language},
CADE-28, Lecture Notes in Computer Science \textbf{12699} (2021), 625--635.

\bibitem{Mathlib}
The mathlib Community,
\emph{The Lean mathematical library},
CPP 2020, 367--381.

\bibitem{Clay}
A.~Jaffe and E.~Witten,
\emph{Quantum Yang--Mills theory},
Clay Mathematics Institute Millennium Prize Problem description, 2000.

\end{thebibliography}

\end{document}
